PHMGA is analyzed as an uneven evidence set rather than as a completed benchmark. The record contains accepted Ottawa rows, rejected RM101 rows, pending backend rows, and a simulated RM101 mechanism signal. That mix creates the condition of interest: a result may carry a strong metric while the admissible claim remains narrow. The analysis starts from evidence units, each defined by the protocol, expected output, source artifact, status, and review or approval state available before interpretation.

Interpretation follows source status before score magnitude. Accepted, rejected, pending, proxy, and mechanism-level sources inherit the limits in Table~\ref{tab:evidence-status-policy}; a prose claim, table entry, or caption can be no stronger than that status permits. When a stronger reading is possible but not yet licensed, it is recorded as a withheld claim together with the evidence that would be needed to make it.

Before a result is used, the manuscript asks four questions: what was measured, which source verdict is attached, what claim that verdict permits, and what evidence would be required to go further. A stream supports the method only when those answers remain visible in the sentence, table, or caption that reports the result. A high metric can then be reported without being used as backend ranking, selection, or performance evidence by itself.

The evaluation then checks three claim-promotion errors. Proxy promotion occurs when a synthetic or offline signal is written as benchmark or real-data evidence. Authority inversion occurs when metric rank overrides source status, for example when a pending or rejected row is used for backend selection. Mechanism inflation occurs when a simulated or mixed-domain signal is written as online or out-of-domain performance. In PHMGA, the synthetic/offline stream tests proxy promotion, the Stage B stream tests authority inversion, and the RM101 stream tests mechanism inflation.

Closure is assigned by the source verdict or approval record, not by the fluency of the paragraph that reports it. For reproducibility, each reported result keeps the identifiers needed to reconstruct provenance: protocol, source record, source artifact, evaluator, result schema, and model/provider metadata when agents are used.
